\chapter{Исходные данные и раздельно проверяемые компоненты метода моделирования ECAP}

\section{Выбор коннектомного графа и установление происхождения исходных данных}
\ResultPending{Будут указаны версия графа, происхождение данных и условия использования.}

\section{Нейронные и поведенческие наблюдения Drosophila и их соответствие выбранному коннектому}
\ResultPending{Будет проверено соответствие наблюдений выбранному коннектому по стадии развития и исследуемой нейронной цепи.}

\section{Выбор доступных записей ECAP, независимой выборки и описание условий регистрации}
\ResultPending{Будут выбраны доступные записи с известными условиями регистрации и возможностью разделения по участникам.}

\section{Коннектомно-ограниченная модель нейродинамики и генерация синтетических данных}
\ResultPending{Будут разработаны модель и генератор с фиксацией воздействий и начальных условий.}

\section{Биофизический оператор рекрутирования волокон, распространения возбуждения и регистрации ECAP}
\ResultPending{Возбуждение волокон и регистрация будут разделены без повторного расчёта рекрутирования.}

\Needspace{7\baselineskip}
\section{Методика подготовки данных и воспроизведения расчётов}
\ResultPending{Будут определены форматы, единицы, происхождение и контрольные суммы данных.}

\section*{Выводы по главе}
\ResultPending{Выводы будут сформулированы после реализации задач главы.}
